%!TEX root = main.tex

\section{Introduction}

% With the rise of the Internet of Things (IoT), the number and diversity of connected devices is on the rise.
% %
% Leaders of the industry, such as Cisco, predict that 50 billion devices will be connected together by 2020~\cite{DaveEvans-IoT}.
% %
% The goal of the IoT is to link\deleted{ this diversity of} these objects and promote machine-to-machine interactions.
% %
% These interactions enable the capture and/or distribution of data that can be turned into information\deleted{ and feedback} delivered to end users.

The advent of the Internet of Things is leading to a revival of interest in mobile ad-hoc networks (MANETs).
%
A MANET is a mobile network that does not rely on a fixed infrastructure.
It is instead a self-configuring network: 
Nodes connect directly to each other, and there is no guarantee for end-to-end connectivity.
%
The topology of a MANET changes frequently due to the mobility of its participants.
%
Nodes are connected using a wireless network, there is no central entity in charge of routing, and most participants operate on battery-powered devices.
%
Applications based on MANETs~\cite{overview_adhoc_networks} range from sensor networks to end-user mobile applications, such as systems for traffic jam prevention or alert dissemination.

Many network functions in MANETs, such as routing, rely on the \textit{broadcast operation} that delivers a message from a source node to all other nodes in the network.
%
Efficiently supporting broadcast in MANETs has received considerable attention.
%
In particular, existing solutions aim to:

\vspace{-1mm}
\begin{itemize}
	\item \textit{Maximize reachability.} 
	This is defined as the percentage of nodes that receive a message. 
	Ideally, every broadcast messages must reach all nodes.
	
	\item \textit{Minimize energy consumption and bandwidth usage.}
	The efficiency of an algorithm is measured by the amount of energy, processing power and network bandwidth it uses.
	
	\item \textit{Be mobility-tolerant and density aware.}
	A broadcast algorithm should guarantee high reachability and cope with mobility patterns under a large variety of scenarios, where the connectivity between nodes is changing constantly.
\end{itemize}

Many algorithms have been proposed to support broadcast in MANETs~\cite{Ovalle-Martinez2006, Qayyum2002, Wu2001, Cartigny:2003:BNR:2727141.2727459, Wu2004, Stojmenovic:2002:DSN:506156.506158}.
%
All these algorithms implement the broadcast with a form of \emph{controlled flooding}: 
Since the transmission range of a source node is typically much smaller than the network deployment area, broadcasting is achieved by message forwarding between nodes.
%
Broadcast algorithms for MANETs mostly differ in the way they control this forwarding operation.
%
Particularly, the control happens in space and time by deciding \emph{which} nodes should forward a message and \emph{when} this operation should take place.
%
Algorithms can also be distinguished based on their use of, or independence from, knowledge about the network topology, nodes locations, and network status.

Broadcast algorithms are largely influenced by external factors such as the density of the network, and the mobility rate of its nodes.
%
The performance of broadcast protocols has been studied individually~\cite{YiGerla2003, Garbinato2010} and extensive work have been done on their classification~\cite{Ruiz2015, Fei1311289}.
%
However, and to the best of our knowledge, there is still no comprehensive experimental comparison of existing solutions, that would allow an understanding of different broadcast algorithms under different deployment scenarios.
%
% There is also a lack of a common set of \textit{broadcasting metrics} followed by the network community to overview the assets of solutions one can find at the literature. %% ER this is a bit too harsh and we do not really solve this problem either. We can safely remove.

\noindent
\textbf{Contributions}.
%
Our goal in this paper is to explore the relative performance and costs of representative MANET broadcast algorithms in a variety of deployment conditions, in order to ease the task of selecting the most appropriate algorithm for a given situation.
%
In addition to simple flooding, we consider four algorithms that are representative of their class:
\begin{itemize}
	\item The \emph{CDS-based} approach builds a source-independent approximation of the minimum Connected Dominating Set (CDS)~\cite{Wu2001}.
	The CDS uses a subset of the nodes as relays, while other nodes are pure receivers.
	
	% \item
	% \changed{\emph{MPR} minimize the number of forwarding nodes building a CDS from the node that initiates a broadcast session~\cite{Qayyum2002}.}
	% \deleted{\emph{MPR} maintains an overlay structure to control the broadcast process. This overlay is built using a heuristic that aims at minimizing the number of forwarding nodes~\cite{Qayyum2002}.}
	
	\item
	\emph{MPR} also uses a subset of nodes as relays but employs heuristics aiming at minimizing their number~\cite{Qayyum2002}.
	
	\item \emph{ABBA} controls the dissemination process in time, by setting up timers that must elapse with no duplicate reception to trigger the forwarding.
	A timer duration is inversely proportional to the probability that the same message will be received from the node neighbors~\cite{Ovalle-Martinez2006}.

	\item \emph{ProbFlood} is a probabilistic approach where the likelihood of forwarding a message is proportional to the measured local density around each node~\cite{Cartigny:2003:BNR:2727141.2727459}.
\end{itemize}

Our study aims at helping practitioners deciding what is the most appropriate algorithm for a given deployment scenario.
%
Our results show that there is no one-size-fits-all to broadcast in MANETs.
%
% We identify which are the algorithms that better match the studied deployment scenarios. %% ER duplicate of beginning of paragraph
%
Our categorization of algorithms can also form the basis of an adaptive broadcast approach that would change its behavior at runtime based on contextual information.
%
% Finally, we outline an adaptive broadcast mechanism in which the algorithm to use is always a good fit for a given deployment scenario. %% ER we do not, 

\raziel{not required in a short paper}

\deleted{
  \noindent
  \textbf{Outline}.
  %
  The rest of this paper is organized as follows.
  %
  Section~\ref{sec:background} presents the algorithms we evaluate in our experiments.
  %
  In Section~\ref{sec:related_work} we survey related work on the evaluation of broadcast algorithms for MANETs and discuss its limitations.
  %
  We present our experimental setup in Section~\ref{sec:scenarios} and our evaluation results in Section~\ref{sec:evaluation}.
  Section~\ref{sec:discussion} summarizes the main findings of this paper.
  %
  Finally, we discuss the perspective of this work in Section~\ref{sec:perspectives} and conclude in Section~\ref{sec:conclusion}.
}

%\er{This is a proposal for an introduction structure, to be discussed}

%\begin{itemize}
%	\item Definition and importance of MANETs
%	\item Characteristics of MANETs
%	\begin{itemize}
%		\item no access to wired network connection
%		\item use a wireless communication medium
%		\item operate on battery
%		\item limited processing power
%		\item mobile in space
%	\end{itemize}
%	\item Operating MANETs rely on a set of fundamental operations, among which broadcast is of great importance \er{ideally we can link to work that relies on broadcast to form routing paths}
%	\item Challenges in MANET networks: limited energy, mobility, wireless communication and collisions, limited knowledge of the topology, etc.
%	\item Broadcast received considerable attention in order to maximize coverage, minimize cost, support mobility patterns, and minimize the consumption of energy
%	\item Protocols in the literature seek to control the flooding of the broadcast message by selecting which node should transmit wirelessly, and when:
%	\begin{itemize}
%		\item Some protocols control the diffusion protocol in space, determining either in advance (topology-based) or during the dissemination, which node is eligible to forward the message when it receives it
%		\item Other protocols control the diffusion protocol in time, determining when a received message should be forwarded and under which conditions, avoiding unnecessary collisions and retransmissions
%	\end{itemize}
%	\item Protocol efficiency depends on their ability to use knowledge about the network topology, node location, and 
%	\item But the efficiency of protocols is also largely influenced by external factors, and primarily by the density of the network, and the mobility rate of its nodes. \er{not sure we should mention other aspects, such as the loosiness of channels, the presence of obstacles, etc.?}
%	\item The relative efficiency of each of the broadcast protocols against the simple and inefficient flooding algorithm, where no control is performed on the flooding operation, has been studied but the relative performance of advanced protocols between them is largely unstudied
%	\item We propose to explore the relative performance and costs of representative MANET broadcast algorithms in a variety of deployment conditions.
%	\item Our study considers three \er{more?} protocols, in addition to the baseline flooding algorithm, that are representative of their class:
%	\begin{itemize}
%		\item ABBA controls the dissemination process in time, by setting up timers prior to dissemination that are inversely proportional to the probability that the message will be received from its neighbors from other sources
%		\item CDS builds and maintain an overlay independently from dissemination, allowing to control the latter in space by selecting which node should or should not propagate incoming messages
%		\item MPR is another well-known algorithms that maintains an overlay structure to control the broadcasting process in space. This overlay is built using a heuristic that aims at minimizing the number of relay nodes required to reach full coverage.
%	\end{itemize}
%	\item Our study allows answering the question of what is the most appropriate protocol for a certain deployment scenario, and concludes that there is no one-size-fits-all in MANET broadcasting \er{provide summary of main findings}
%\end{itemize}
